\documentclass[11pt]{amsart}
\usepackage[T1]{fontenc}
\usepackage[utf8]{inputenc}
\usepackage{amssymb}
\usepackage[margin=1.1in]{geometry}
\usepackage[hidelinks]{hyperref}
\newtheorem{theorem}{Theorem}
\newtheorem{proposition}[theorem]{Proposition}
\newtheorem{lemma}[theorem]{Lemma}
\newtheorem{corollary}[theorem]{Corollary}
\theoremstyle{definition}
\newtheorem{definition}[theorem]{Definition}
\newtheorem{axiom}[theorem]{Axiom}
\theoremstyle{remark}
\newtheorem{remark}[theorem]{Remark}
\newcommand{\Om}{\Omega}
\newcommand{\Map}{\mathrm{Map}_*}
\newcommand{\Z}{\mathbb{Z}}
\newcommand{\code}[1]{\texttt{#1}}
\newcommand{\Hom}{\mathrm{Hom}}
\newcommand{\hc}[2]{[#1,#2]_*}
\newcommand{\Fib}{\mathrm{Fib}}
\newcommand{\tr}[1]{\|#1\|}
\newcommand{\cls}[1]{[#1]}
\newcommand{\kb}{\bar\kappa}
\newcommand{\Fh}{\hat F}
\newcommand{\cj}{\mathrm{conj}}
\newcommand{\ucj}{\mathrm{unconj}}
\newcommand{\Unl}{\mathsf{Unloop}}
\title{$\pi_4(S^3)=\Z/2$ by a Postnikov universal property:\\ the main ideas of a formal proof}
\date{28 October 2026}
\begin{document}
\maketitle

\noindent\textbf{Attribution.} David W\"arn developed the
universal-property strategy used here (one identifies a candidate Postnikov stage by a
\emph{marked universal property} and compares it with the truncated sphere on homotopy
classes), together with the delooping constructions (Eilenberg--MacLane spaces as
iterated torsors, and the criterion for a map of loop spaces to be the looping of a
map~\cite{DW1}) and the cup-product construction. Our contribution was to develop the
argument in detail and formalise it in Unit; the formal proof is the development
\code{sag/pi3-s2-postnikov}, and the detailed notes are listed in \S\ref{s:formal}.
Owen Milner formulated the S\'{\i}nh invariant in homotopy type theory.

\smallskip\noindent\textbf{The result.} In univalent type theory with replacement,
and assuming only the universal property of $S^3$ (Axiom~\ref{a:s3}),
\[\pi_4(S^3)\cong\Z/2\quad\text{as groups}\]
(\code{Pi4Sphere.Pi4S3.pi4-S3}, Theorem~\ref{t:final}). No additional higher-inductive
construction is assumed beyond the stated universal property of the sphere; no pushout,
no smash product and no spectral sequence is used; the two homotopy-theoretic
inputs are the cohomology computations (H3), (H4) of $K(\Z,2)$, proved here from the
circle and the cup square.

\section{Setting}

The foundations are those of \code{sag/em}: univalence and Rijke's replacement
(images and their surjectivity) as parameters; function extensionality,
propositional truncation, set quotients and the finite truncations $\tr X_m$ are
derived. Eilenberg--MacLane spaces $K(A,n)$ of an abelian group $A$ are built as in
\cite{DW1}: $K(A,1)$ is the type of $A$-torsors and $K(A,n{+}1)$ is the type of
$K(A,n)$-torsors, i.e.\ of merely inhabited types with a homogeneous structure. We
write $e_n:\Om K(A,n{+}1)\simeq_*K(A,n)$, $K(h,n)$ for the map induced by a
homomorphism $h$, and $F;G$ for composition of pointed maps (first $F$).

For pointed $X,Y$, $\hc XY$ is the \emph{set} of pointed homotopy classes, the set
truncation of $\Map(X,Y)$; classes compose and form a category. Two classes are equal
when their representatives are merely equal; no representative is ever chosen.
Homotopy groups are $\pi_3X=\pi_1(\Om^2X)$ and, by definition,
\[\pi_4X=\pi_3(\Om X)=\pi_1(\Om^3X),\]
functorial on maps and on classes (looping descends to classes since $\hc{\Om X}{\Om Y}$
is a set).

\begin{axiom}[$S^3$]\label{a:s3}
A type $S^3$ with a point $\mathsf{base}$ and $\mathsf{surf}:\Om^3(S^3,\mathsf{base})$
such that for every $X$, $f\mapsto(f\,\mathsf{base},\mathsf{ap}^3_f\,\mathsf{surf})$ is
an equivalence $(S^3\to X)\simeq(x:X)\times\Om^3(X,x)$.
\end{axiom}

It is stated as the parameters of a module, not as a global assumption. Its pointed
form is $\Map(S^3,X)\simeq\Om^3X$, $F\mapsto(\Om^3F)(\mathsf{surf})$, natural in $X$.

\section{The strategy: a marked universal property}

The proof is a Yoneda argument on classes; the same argument, one level down, gives
$\pi_3(S^2)=\Z$ (formalised first, \code{doc/pi3-s2-postnikov.pdf}).

Let $T=\tr{S^3}_4$ with unit $\eta:S^3\to_*T$. From the axiom, $S^3$ is 2-connected
(into a 2-type $Y$ evaluation at the base point is an equivalence, so $\tr{S^3}_2$ is
contractible), hence $T$ is a 2-connected 4-type, $\eta^*:\Map(T,Y)\to\Map(S^3,Y)$ is an
equivalence of pointed mapping types for every 4-type $Y$, hence a bijection
$\hc TY\to\hc{S^3}Y$ on homotopy classes, and $\pi_4\eta$ is an isomorphism (the fibres of $\Om^k\eta$ are
$(4-k)$-connected).

\begin{proposition}[two instances suffice]\label{p:two}
Let $P$ be a 4-type and $s\in\hc{S^3}P$ such that $s^*:\hc PY\to\hc{S^3}Y$ is a
bijection for $Y=P$ and for $Y=T$. Then $\pi_4S^3\cong\pi_4T\cong\pi_4P$.
\end{proposition}
\begin{proof}
Let $u\in\hc TP$ with $\cls\eta;u=s$ ($\eta^*$ at $P$) and $t\in\hc PT$ with
$s;t=\cls\eta$ ($s^*$ at $T$). Then $s;(t;u)=\cls\eta;u=s=s;\mathrm{id}$, so
$t;u=\mathrm{id}_P$; and $\cls\eta;(u;t)=s;t=\cls\eta$, so $u;t=\mathrm{id}_T$. Hence
$\pi_4u$ and $\pi_4t$ are inverse isomorphisms.
\end{proof}

So the whole proof consists in constructing a 4-type $P$ with $\pi_4P=\Z/2$ and a
class $s$ with this universal property. The candidate is the delooping $BR$ of
$R=\Fib\kb$, the fibre of the cup-square $k$-invariant $\kappa:K(\Z,2)\to_*K(\Z,4)$ with
its coefficient reduced modulo~2, $\kb=\kappa;K(\rho,4)$ (\S3). No result about $S^2$ is
used.

\section{The candidate $BR$}

\noindent\textbf{The cup square.} Following W\"arn's construction, the circle property
of $K(\Z,1)$ and delooping give a bipointed map $c_2:K(\Z,2)\to_*\Map(K(\Z,2),K(\Z,4))$,
normalised so that looping it and evaluating at the generator of $\pi_1K(\Z,1)$ gives
back the identity of $K(\Z,2)$; the cup square is its diagonal, $\kappa(x)=c_2(x)(x)$
(\code{doc/cup-square.pdf}).

\noindent\textbf{Cohomology of $K(\Z,2)$.} For every abelian group $B$:
\begin{itemize}
\item[(H3)] $\hc{K(\Z,2)}{K(B,3)}$ is a point (from the circle property;
  \code{doc/circle-torsors.pdf});
\item[(H4)] $\Phi_B:\Hom(\Z,B)\to\hc{K(\Z,2)}{K(B,4)}$,
  $h\mapsto\cls{\kappa;K(h,4)}$, is a bijection (\code{doc/h4.pdf}): every class is the
  diagonal of a unique bipointed map, and bipointed maps are classified by
  $g\mapsto c_2;K(g,4)_*$.
\end{itemize}
With $P_2=\Fib\kappa$ this gives the mapping theorem and $\pi_3(S^2)=\Z$.

\noindent\textbf{Reduction mod 2.} Let $\rho:\Z\to\Z/2$, $\kb=\kappa;K(\rho,4):K(\Z,2)\to_*
K(\Z/2,4)$ and $R=\Fib\kb$, a simply connected 3-type with $\pi_2R=\Z$ and
$\pi_3R\cong\Z/2$ as groups.

\begin{theorem}[$R$ deloops; \code{doc/cross-term.pdf}]\label{t:deloop}
There is $\Fh:K(\Z,3)\to_*K(\Z/2,5)$ with $\Om\Fh=e_2;\kb;e_4^{-1}$. Hence
$BR=\Fib\Fh$ is a pointed 2-connected 4-type with $\Om BR\simeq_*R$.
\end{theorem}
\begin{proof}[Idea]
By W\"arn's delooping criterion (\cite{DW1}, Lemma~9), $f=e_2;\kb;e_4^{-1}$ is a looping
as soon as, for every loop $p$, $f$ agrees with its translate $q\mapsto(fp)^{-1}\cdot f(p\cdot q)$;
at these levels the space of such data is a proposition, so no coherence for the
pointing in $p$ is needed. The defect is the \emph{cross term}
``$\kb(x+y)-\kb(x)-\kb(y)$'', a bipointed map; by the classification of (H4) it
vanishes as soon as its diagonal ``$\kb(2x)-2\kb(x)$'' does, and both terms of the
diagonal vanish by construction because $2=0$ in $\Z/2$: doubling on
$\Om K(\Z/2,m{+}1)$ is $K(2,m)=K(0,m)=0$, and precomposing $c_2$ with $K(2,2)$ doubles
its invariant.
\end{proof}

\section{The mapping theorem for $BR$}

A \emph{level-3 triple} is $(A,B,K)$ with $K:K(A,3)\to_*K(B,5)$; its fibre $\Fib K$ is
a pointed 2-connected 4-type, with projection $q:\Fib K\to_*K(A,3)$.

\begin{theorem}[lifting]\label{t:lift}
For a pointed $Q$, suppose (E) every $F:Q\to_*K(A,3)$ has merely $F;K=0$, and (U) every
$d:Q\to_*\Om K(B,5)$ is merely $0$. Then $q_*:\hc Q{\Fib K}\to\hc Q{K(A,3)}$ is a bijection.
\end{theorem}

This is the lifting lemma of the mapping theorem for $\Fib\kappa$
(\code{doc/mapping-fib.pdf}) one level up. It is applied to the two sources $Q=S^3$ and
$Q=BR$. For $S^3$ both conditions are immediate (pointed maps out of $S^3$ are 3-loops,
and $\Om^3K(B,5)$, $\Om^4K(B,5)$ are connected).

\noindent\textbf{(U) for $BR$.} Looping is an equivalence
$\Map(BR,\Om K(B,5))\simeq\Map(R,\Om^2K(B,5))$, and $\hc R{K(B,3)}=0$: a map out of
$R=\Fib\kb$ restricts on the fibre to a map that deloops to some $K(g,4)$, the fibre
sequence forces $\cls{\kb;K(g,4)}=0$, and
\[\Psi_B:\Hom(\Z/2,B)\to\hc{K(\Z,2)}{K(B,4)},\qquad g\mapsto\cls{\kb;K(g,4)}=\Phi_B(\rho;g)\]
is injective by (H4); then (H3) kills what remains (\code{doc/fibre-h3.pdf}, with a fibre
coefficient group).

\noindent\textbf{(E) for $BR$: the decisive point.} A map $F:BR\to_*K(A,3)$ factors merely
through $p:BR\to_*K(\Z,3)$ as $F=p;G$ with $G:K(\Z,3)\to_*K(A,3)$ (the fibres of $p$ are
3-connected). Let $M=G;K:K(\Z,3)\to_*K(B,5)$; then merely $F;K=p;M$, and the obstruction
is the class of $H=\ucj(\Om M):K(\Z,2)\to_*K(B,4)$, i.e.\ $h\in\Hom(\Z,B)$ with
$\Phi_B(h)=\cls H$. The fibre sequence of $BR$ kills it only if $h$ factors through
$\rho$, i.e.\ if $2h=0$. So one must show that \emph{for every target the obstruction is
2-torsion}. The point is that $H$ is the looping of a map \emph{by construction}:

\begin{lemma}[necessity]\label{l:N}
If $\Phi_B(h)=\cls{\ucj(\Om\hat G)}$ for some $\hat G:K(\Z,3)\to_*K(B,5)$, then $2h=0$.
\end{lemma}
\begin{proof}[Idea]
A looping commutes with squaring of loops, and squaring is $K(2,2)$ on the source and
$K(2_B,4)$ on the target; so $\cls{K(2,2)};\cls H=\cls H;\cls{K(2_B,4)}$. A computation
with the bipointed classification gives $(2x)^2=4x^2$, i.e.\
$\cls{K(2,2)};\cls\kappa=\cls\kappa;\cls{K(4,4)}$; its only new ingredient is the
naturality of the invariant of a bipointed map under precomposition in the second
variable. Hence $\Phi_B(4h)=\Phi_B(2h)$, and $2h=0$ since $\Phi_B$ is injective.
\end{proof}

\noindent This is the central reason why $\Z/2$ appears: looping forces doubling to act
\emph{linearly} on the obstruction ($h\mapsto2h$), whereas the cup square makes it act
\emph{quadratically} ($h\mapsto4h$); comparing the two forces $2h=0$.

Then $h=\rho;g$, so merely $H=\kb;K(g,4)$, $\Om M=\Om(\Fh;K(g,5))$, and since looping is
injective on maps $K(\Z,3)\to_*K(B,5)$, merely $M=\Fh;K(g,5)$ and
$F;K=p;M=(p;\Fh);K(g,5)=0$. This is why the target is presented one level up, with a
$k$-invariant $K(A,3)\to K(B,5)$: every obstruction is then a looping, and no coherence
cell has to be normalised.

\noindent\textbf{The marked class.} From the axiom and the circle property one builds
$\sigma_3:S^3\to_*K(\Z,3)$ with $(\Om^3\sigma_3)(\mathsf{surf})$ the generator; $\sigma_3^*$
is an equivalence on maps into 3-types (a composite of five equivalences). Let
\[s=p_*^{-1}\cls{\sigma_3}\in\hc{S^3}{BR}\]
(Theorem~\ref{t:lift} at $Q=S^3$ for the triple $(\Z,\Z/2,\Fh)$). Then $s^*$ is a
bijection into every $K(A,3)$, and the commuting square of $q_*$ for the two sources gives:

\begin{theorem}[mapping theorem]\label{t:main}
For every level-3 triple, $s^*:\hc{BR}{\Fib K}\to\hc{S^3}{\Fib K}$ is a bijection; hence
so is $s^*:\hc{BR}Y\to\hc{S^3}Y$ for every $Y$ merely pointed equivalent to some
$\Fib K$.
\end{theorem}

\section{Presentations of 2-connected 4-types}

\begin{theorem}\label{t:pres4}
Every pointed 2-connected 4-type $Y$ is merely pointed equivalent to $\Fib K$ for a
level-3 triple $(A,B,K)$.
\end{theorem}
\begin{proof}[Idea]
As the presentation of $\tr{S^2}_3$ (\code{doc/postnikov-presentation.pdf}), one level up.
Let $V=\tr Y_3$; it is $K(A,3)$ with $A=\Om^3V$ (a 2-connected 3-type is recognised
by delooping its loop space twice). The fibres $E_z$ of $Y\to V$ are 3-connected
4-types, hence $K(B,4)$ with $B=\Om^3(\Om E_{v_0})$. The structures $K(B,4)\simeq_*E_z$
form a set over $Y$, and a section through the given one exists since $Y$ is
2-connected; so every $E_z$ carries a homogeneous structure, i.e.\ is a
$K(B,4)$-torsor, and $z\mapsto E_z$ is a map $k':V\to K(B,5)$ with $\Fib k'\simeq Y$.
Take $K=e_V;k'$; the pointing is only mere since $Y$ is connected.
\end{proof}

The group $B$ is not identified: for $Y=T$ it is $\pi_4T$, the unknown. Since the goal
``$s^*$ is a bijection'' is a proposition, the mere presentation of $T$ suffices.

\section{Conclusion}

\begin{theorem}\label{t:final}
Under Axiom~\ref{a:s3}, $\pi_4(S^3)\cong\Z/2$ as groups.
\end{theorem}
\begin{proof}
Theorem~\ref{t:main} applies to $Y=BR=\Fib\Fh$ directly and to $Y=T$ through
Theorem~\ref{t:pres4}. Proposition~\ref{p:two} with $P=BR$ gives
\[\pi_4S^3\cong\pi_4T\cong\pi_4BR=\pi_3(\Om BR)\cong\pi_3R\cong\Z/2.\qedhere\]
\end{proof}

\section{The formalisation}\label{s:formal}

The proof is checked in Unit, in \code{sag/pi3-s2-postnikov} (\code{./segc
sag/pi3-s2-postnikov all}: 83 modules; a fresh check, 4038 declarations, 0 failed, about
2\,s). The final declaration is
\begin{center}
\code{Pi4Sphere.Pi4S3[S3, base, surf, ump3].pi4-S3 : GIso[pi4 S3 base, ab-grp Z2ab]}
\end{center}
whose only parameters beyond the foundations (univalence, images and their surjectivity)
are the $S^3$ axiom; the 4-truncation is constructed inside. The $S^2$ axiom is not used.

\smallskip
\begin{tabular}{@{}ll@{}}
\code{CupSquare}, \code{H3Proof}, \code{H4Proof} & the cup square, (H3), (H4)\\
\code{CrossTerm}, \code{DeloopR} & the cross term vanishes; $\Fh$, $BR$, $\Om BR\simeq R$\\
\code{Necessity}, \code{FibreH3} & Lemma~\ref{l:N}, $(2x)^2=4x^2$; $\hc R{K(B,3)}=0$\\
\code{MapBR} & Theorem~\ref{t:lift}; (E) and (U) for $BR$\\
\code{Sphere3Axiom}, \code{SourceS3} & the axiom; $\sigma_3$; $s$; Theorem~\ref{t:main}\\
\code{Sphere3Trunc} & $S^3$ is 2-connected; $T=\tr{S^3}_4$; $\pi_4$ on classes\\
\code{Presentation4} & Theorem~\ref{t:pres4}; the mapping theorem at $T$\\
\code{Pi4Sphere} & Proposition~\ref{p:two}; Theorem~\ref{t:final}
\end{tabular}

\smallskip
For every declaration the checker reports which definitions it had to unfold; the
\code{.opens} files of the development record, declaration by declaration, the
definitions allowed to be unfolded, and the check against them passes for every
declaration. These specifications were generated from the checker's report and reviewed
by kind (which kinds of unfolding each text step justifies); the textual audit, entry by
entry against the notes, is not yet complete. The later parts (the cup square, (H3),
(H4), the fibre lemma, the mapping theorem for $\Fib\kappa$, the presentation of
$\tr{S^2}_3$, and everything from the cross term onwards) were written from an
informal note first; the earliest modules (the sphere axiom, homotopy classes and
groups, the realisation of Postnikov triples, finite truncations) were written before
their notes. Detailed notes:
\code{br-universal.pdf} (this proof in full), \code{cross-term.pdf},
\code{cup-square.pdf}, \code{h4.pdf}, \code{fibre-h3.pdf}, \code{circle-torsors.pdf},
\code{mapping-fib.pdf}, \code{postnikov-presentation.pdf}, \code{finite-truncation.pdf},
and \code{pi3-s2-postnikov.pdf} for $\pi_3(S^2)=\Z$.

\begin{thebibliography}{9}
\bibitem{DW1} D.~W\"arn. \emph{Eilenberg--MacLane spaces and stabilisation in homotopy
  type theory}. arXiv:2301.03685, 2023.
\end{thebibliography}
\end{document}
